% master.tex : master-fil for projektet
% ------------------------------------------------------------------------------
% Dette er hovedfilen for projektet, hvori indhold fra alle input-filer (tekst,
% billeder, litteraturdatabaser, osv.) samles

% Dokumenttypen 'book' er valgt pga. dens mange fleksible indstillinger
% Se https://tex.stackexchange.com/a/36989/118167
\documentclass[11pt,a4paper,twoside,openright,danish]{book}

% Variabler, som bruges til automatisk at indsætte titel, forfattere, osv. på
% forsiden og titelbladet.
\def \projecttitle       {Optimering}
\def \projectsubtitle    {Lineær programmering}
\def \projecttheme       {Lineær algebra}
\def \projectdegree      {Matematik}
\def \projectperiod      {Efterårssemesteret 2020}
\def \projectnumber      {P2}
\def \projectgroup       {B303c}
\def \projectauthors     {
	Jens Koldkur\\
	Jonas Bach Christensen\\
	Julie Havbo Lund\\
	Mads Bjerregaard Kjær\\
	Mathias Gammelgaard\\
	Rasmus Jespersgaard
}
\def \projectsupervisors {
  Horia Cornean
}

% Preamblet indeholder alle de indstillinger og makroer, som skal indsættes for
% hovedindholdet, og i denne skabelon samles det i filen aaumath.sty, som
% definerer en pakke, der kan indlæses med \usepackage.
\usepackage{aaumath}

% Dokumentets indhold indsættes mellem \begin- og \end-makroerne for
% 'document'-blokken
\begin{document}

% Dokumentets 'front matter' tælles ikke med ifm. antal sider og nummereres med
% romerske tal. Herunder hører f.eks. forsiden, titelbladet, forordet og
% indholdsfortegnelsen.
\frontmatter
% incl/misc/frontpage.tex : rapportens forside
% ------------------------------------------------------------------------------


\backgroundsetup{
  scale = 1,
  angle=0,
  opacity=1,
  contents = {
    \includegraphics[width=\paperwidth,height=\paperheight]{fig/img/aau/waves.pdf}
  }
}
\BgThispage
\pdfbookmark[0]{Forside}{forside}
\begin{titlepage}
  \centering
  \phantom{}
  \vspace{2cm}

  % AAU-segl
  \begin{minipage}[c]{0.2\paperwidth}
    \centering
    \makebox[0pt]{
      % fig/tikz/aau-badge.tex : AAU-logo til forsiden
% ------------------------------------------------------------------------------

\begin{tikzpicture}
  % Tegn hvid cirkel og tilføj det gennemsigtige, blå logo ovenpå
  \node[circle,color=white,fill=white,minimum size=1.175\textwidth] at (0,0) {};
  \node at (0,0) {\includegraphics[width=\textwidth]{fig/img/aau/logo-circle.pdf}};
\end{tikzpicture}

    }
  \end{minipage}

  % Hovedindhold
  \vspace{4cm}
  {\fontfamily{bch}\selectfont
    \fboxsep0pt\colorbox{white}{
      \begin{minipage}{\textwidth}
        \centering
        \color{AAUblue1}

        \vspace{2em}
        {\Huge\bfseries\projecttitle}

        {\Large\bfseries\projectsubtitle}

        \bigskip
        \parbox{\textwidth}{\centering\large\projectauthors}

        \bigskip
        {\projectdegree}
        \vspace{2em}
      \end{minipage}
    }
  }

\end{titlepage}

% incl/misc/titlepage.tex : rapportens titelblad
% ------------------------------------------------------------------------------
% Titelbladet genereres af makroen \aautitlepage, som er defineret i
% /incl/pre/ext/aautitlepage.sty


\pdfbookmark[0]{Titelblad}{titelblad}
\aautitlepage{
  \projectinfo{
    \projecttitle
  }{
    \projecttheme
  }{
    \projectperiod
  }{
    Gruppe \projectgroup
  }{
    \parbox[t]{\textwidth}{\projectauthors}
  }{
    \parbox[t]{\textwidth}{\projectsupervisors}
  }{
    \today
  }
}{
  \textbf{Institut for Matematiske Fag}\\
  Skjernvej 4A\\
  DK-9220 Aalborg Ø\\
  \href{http://math.aau.dk}{http://math.aau.dk}
}{
  % incl/misc/abstract.tex : projektets abstract
% ------------------------------------------------------------------------------
% Et abstract er et kort resume af rapporten, som vises på titelbladet
This paper contains an examination on linear programming problems and their geometrical solutions in relation to the simplex method. 
The paper contains a theoretical introduction to linear algebra together with important definitions and theorems related to this field of study. 
Afterwards the solution of linear programming problems and notations are introduced, which leads into the geometrical representation of  linear programming problems through polyhedrons. Here it is proven that an optimal solutions exists in a corner of such a geometrical figure.
Lastly the report provides a link between the use of the simplex method, which is a solution algorithm for linear programming problems, and the geometrical representation, aswell as the practical application of this.
Different implementations of the simplex method are illustrated her aswell and discussed in regards to their time complexity. 
}

% incl/misc/contents.tex :
% ------------------------------------------------------------------------------

\pdfbookmark[0]{Indhold}{indhold}

% Indstillingerne i denne fil er grupperet, så de ikke påvirker andre dele
\begingroup

% Slå twoside fra midlertidigt for at undgår sideskift
\makeatletter
\@twosidefalse
\makeatother

% Placer indholdsfortegnelsen på sin egen side (bedst når den kun fylder en side)
\tableofcontents
\clearpage

% Placer oversigtslisterne på efterfølgende sider
%\let\clearpage\relax
%\listoffigures
%\listoftables
%\listofalgorithms
%\lstlistoflistings

\endgroup


% Dokumentets 'main matter' (hovedindhold) er der, hvor det meste indhold skal
% sættes ind. Sider og overskrifter er nummererede med arabiske tal.
\mainmatter

% Input-filer bør opdeles således, at hver fil svarer til et kapitel. Makroen
% \include indsætter et sideskift og indholdet fra den givne stil.

% \include{incl/main/example1}
% \include{incl/main/example2}
% ...


% INDLEDNING
	\chapter{Forord}

Jens Koldkur,
Jonas Bach Christensen,
Julie Havbo Lund,
Mads Bjerregaard Kjær,
Mathias Gammelgaard og
Rasmus Jespersgaard.
	
2. semester 
	
bacheloruddannelsen Matematik på Aalborg Universitet
	
2020
	
udarbejdet fra 3 februar til 27 maj 2020
	
kravende 
	
temaet 
	
Rapporten er nummereret i afsnit, definitioner, figurer m.m.
Beviser afsluttes med $\square$, som indikerer, at beviset er gennemført. \\\\
% \square skulle have været \qed!

En tak til vores vejleder Horia Cornean
	\chapter{Indledning}
%To til tre sider 
%Vi skal argumentere det er interessant at kigge på lineær programmering
%Kort beskrivelse af hvad vi har lavet 
%Eventuelt 
%   Kapitel 1: bla bla 
%   Kapitel 2: bla bla 
%   m.m.
Optimeringsproblemer er en fundamental del af matematikken, og i den anvendte matematik udgøre disse en vægtig del af matematikkeres virke.
I denne rapport vil liniære optimeringsproblemer undersøges. Dele af forskningsfeltet betragter disse som havende ophav ved Danzig og dennes udarbejdelse af simplexmetoden(ref) i 1947.
Der er dog tilstadighed diskussion om hvor liniære programeringsproblemer som matematisk felt stammer fra og andre peger på Sovjetunionen og behovet for at optimere produktionen i en planøkonomisk kontekst som ophav.(ref2)  
Forudsætningen for at Danzig arbejdede med liniære programeringsproblemer var hans ansættelse ved Royal Airforc(ref) og en applikation blev i forbindelse med den vestligt kontrolerede sektor af Berlin.
Vestberlin lå dybt inde i den sovjetisk kontrolerede sektor af Tyskland og som følge af spændingerne mellem vesteuropa og Sovjet besluttede Josef Stalin at lukke for forsyninger via jorden til Berlin.
Det blev derfor nødvendigt for vestmagterne at etablere en luftbro til Berlin der gjorde det muligt trods de sovjetiske restriktioner at indføre basale fornødehender til de $2,5$ milloner mennesker der levede i byen.
Da en sådan måde at indføre vare på er ressourceintensiv, blev der derfor behov for at gøre dette mest effektivt underlagt de logistiske restriktioner.
Det er derfor et matematisk felt der i sin brug er nært bundet op på praktiske problemstillinger dette værende økonomiske og logistiske.
Det er derfor også et felt der har været nært knyttet op til militære problemstillinger som Danzigs ansættelsesforhold også antyder.
Danzigs bidrag til forskningsfeltet, altså simplexmetoden, er dog stadigt en fundamental del af hvordan løsningen af liniære programmeringsproblemer behandles i dag.
Projektets primære arbejdsområde vil derfor omhandle simplex-metoden samt den geometriske fremstilling af liniære programeringsproblemer.
Det skal dog nævnes at der ligeledes findes optimeringsproblemer der ikke kan beskrives igennem liniære sammenhænge, en afgrænsning bliver her at disse ikke vil beskrives i dette indeværende projekt.
Problemformuleringen bliver derfor som følger:
Problemformulering:
Hvordan kan liniære porgrameringsproblemer anskues geometrisk og hvordan kan disse løses ved hjælp af simplexmetoden?
Med henblik på at besvare denne vil følgende emneområder bliver inddraget:
\begin{itemize}
\item En introduktion til liniær algebra: Afsnittet har til hensigt at indføre passende notation, samt introducere væsentlige metoder og begreber til løsning af liniære programeringsproblemer.
\item Liniære programeringsproblemer: Projektets primære emneområde. Det vil her introduceres hvordan disse fremstilles matematiske samt interessante egenskaber herved.
\item Geometrisk fremstilling: her skal der stå et eller andet når vi ved hvad der er med.
\item Simplexmetoden: Beskrivelse og bevis af Simplex som løsningsalgoritme.
\end{itemize}
	
% LINEÆRE ALGEBRA 
	\chapter{Lineær algebra}
\section{Matricer}
% God arbejdslyst 

	\section{Lineære ligningssystemer }
%

\textbf{En lineær ligning} med variablerne, $x_1,x_2,\ldots ,x_n$, er en ligning som har formen
$$ a_1x_1+a_2x_2+\cdots+a_nx_n=b$$ 
hvor $a_1,a_2,\ldots,a_n$ er reelle tal. 

\begin{defn}{}{}
Et lineært ligningssystem er en mængde af $m$ lineære ligninger  med de samme $n$ variabler, hvor $m$ og $n$ er positive heltal.
Dette kan skrives på formen
\begin{align*}
a_{11}x_1+a_{12}x_2+&\cdots+a_{1n}x_n=b_1\\
a_{21}x_1+a_{22}x_2+&\cdots+a_{2n}x_n=b_2\\
&\vdots\\
a_{m1}x_1+a_{m2}x_2+&\cdots +a_{mn}x_n=b_m
\end{align*}
hvor er $a_{ij}$ betegner koefficienten af $x_j$ i ligning $i$
\end{defn}
\noindent
Løsningen på et sådanne system er en vektor$\begin{bmatrix}
s_1\\
s_2\\
\vdots\\
s_n
\end{bmatrix}
$i $\R ^n$ således at alle ligninger er opfyldt når $x_i$ bliver erstattet med $s_i$.
\begin{eks}\label{eks}
\begin{align*}
x_1-x_3-2x_4-8x_5&=-3 \\
-2x_1+x_3+2x_4+9x_5&=5 \\
3x_1-2x_3-3x_4-15x_5&=-9
\end{align*}
Løsningen på dette system ses i eksempel \ref{eks_gauss}
\end{eks}
\phantom{g}\\\\
Alle lineære ligningssystemer har enten ingen, præcis en eller uendelig mange løsninger.
Alle ligningerne i systemet repræsenteres af linje hvis systemet er i 2 variable, og et plan i 3 variable, osv.
Hvis disse ligningers er parallelle og forskellige, har systemet ingen løsning.
Hvis ligningerne er forskellige, så har systemet præcis en løsning. 
Hvis ligningerne er ens har systemet uendelig mange løsninger.\\
Et system der har en eller flere løsninger kaldes \textit{konsistent}, ellers kaldes det for \textit{inkonsistent}, da der ingen løsninger eksisterer.
\subsection{Ligningssystemer og matricer}
Med udgangspunktet i eksmpel \ref{eks}, kan et lineært ligningssystem opstilles som en matrix ligning A\textbf{x}=\textbf{b}, hvor\\
\begin{center}
$A=
\begin{bmatrix}
1 & 0 & -1 & -2 & -8\\
-2 & 0 & 1 & 2 & 9\\
3 & 0 & -2 & -3 & -15
\end{bmatrix}
$, 
$\textbf{x}=
\begin{bmatrix}
x_1\\
x_2\\
x_3\\
x_4\\
x_5
\end{bmatrix}
$
og
$
\textbf{b}=\begin{bmatrix}
-3\\
5\\
-9
\end{bmatrix}.
$
\end{center}
A kaldes for koefficientmatricen. 
%efter hårdt sammenleje bliver A matricen og b til totalmatricen
Totalmatrixen er dannet ved at kombinere koefficientmatricen med vektoren \textbf{b}.
Totalmatrixen indeholder alle informationer der skal bruges til at løse et lineært ligningssystem.
Ligningssystemet fra eksempel \ref{eks} skrevet som totalmatrix
\begin{equation*}
  A=
\begin{blockarray}{ccccccc}
x_1 & x_2 & x_3 & x_4 & x_5 & b \\
\begin{block}{[ccccc|c]c}
  1 & 0 & -1 & -2 & -8 & -3 \\
  -2 & 0 & 1 & 2 & 9 & 5 \\
  3 & 0 & -2 & -3 & -15 & -9 \\
\end{block}
\end{blockarray}.
\end{equation*}

	\section{Gauss elimination}
%sætning 1.4 kan vi bevise den
En algoritme til løsningen af linære ligningssystemer er  rækkereduktionsalgoritmen, som omdanner en matrix fra en totalmatrice til reduceret trappeform.
Når den bruges til at løse linære ligningssystemer, kaldes denne process for Gauss elimination.
I algoritmen bruges de elemæntere rækkeoperationer for at omdanne en matrix $A$ til trappeform. 
Herefter benyttes rækkeoperationerne igen til at omdanne matricen på trappeform til reduceret trappeform.
Algortimen gennemgår således følgende skridt:
%
\begin{enumerate}
\item Find første ikke-nul søjlse i $A$ fra venstre.
\item Ved rækkeombytning placeres en ikke-nul indgang øverst i pivot-søjlen.
\item Skab nuller under pivot-indgangen øverst i pivot-søjlen ved hjælp af rækkeudskiftning.
\item Øverste række markeres som afsluttet og trin $1-3$ gennemføres nu på den næste række.
\item Alle rækker med pivot-indgange skaleres så alle pivot-indgange er lig $1$.
\item Ved rækkeudskiftning sikres nu $0$er over og under pivot-indgangene.
\end{enumerate}
% Eksempel
Givet ligningssystemet:
\begin{align*}
x_1-x_3-2x_4-8x_5&=-3 \\
-2x_1+x_3+2x_4+9x_5&=5 \\
3x_1-2x_3-3x_4-15x_5&=-9 \\
\end{align*}
\[
\begin{blockarray}{cccccc}
a & b & c & d & e \\
\begin{block}{[ccccc]c}
  1 & 1 & 1 & 1 & 1 & f \\
  0 & 1 & 0 & 0 & 1 & g \\
  0 & 0 & 1 & 0 & 1 & h \\
  0 & 0 & 0 & 1 & 1 & i \\
  0 & 0 & 0 & 0 & 1 & j \\
\end{block}
\end{blockarray}
 \]

\end{document}

\begin{thm}{}{konsitens}
\end{thm}

% Appendicer indsættes inde i en appendices-blok og bliver nummereret med
% bogstaver i stedet for tal
\begin{appendices}
  % \include{incl/app/appendix1}
  % \include{incl/app/appendix2}
  % ..
\end{appendices}

% Dokumentets 'back matter' er til ekstra ting som f.eks. litteraturlisten.
% Overskrifter bliver ikke nummereret her.
\backmatter

% Automatisk litteraturliste baseret på, hvilke kilder, der er blevet refereret
% til i løbet af rapporten.
\bibliographystyle{apalike}
\bibliography{
  incl/bib/books,
  incl/bib/articles,
  incl/bib/software
}

\end{document}
